%% CHAPTER 4: IMPLEMENTATION
\section{CHAPTER 4: IMPLEMENTATION}

\subsection{4.1. Data Collection Module (\texttt{download\_data.py})}

\subsubsection{4.1.1. Overview and Motivation}
The data collection module is the entry point of the X-AURUM pipeline. Rather than relying on a manually curated or periodically downloaded dataset, the Python script is designed to be re-run at any time to refresh the training data with the most recent market information. This automated harvesting approach ensures the model can be periodically retrained on up-to-date price history as gold market conditions evolve.

The script uses the \texttt{requests} library for HTTP communication and the standard \texttt{csv} module for writing the output file. No additional third-party dependencies are required beyond a standard Python installation.

\subsubsection{4.1.2. Binance API Call Structure}
The core API call requests the 1,000 most recent daily candlesticks for PAXG/USDT:

\begin{lstlisting}[language=Python]
import requests, csv

url = "https://api.binance.com/api/v3/klines"
params = {"symbol": "PAXGUSDT", "interval": "1d", "limit": 1000}
response = requests.get(url, params=params, timeout=15)
response.raise_for_status()
klines = response.json()
\end{lstlisting}

The Binance API returns each candlestick as a 12-element list. The elements at indices 1--6 correspond to Open, High, Low, Close, Volume, and Close Time respectively. The script writes these to a CSV file with a descriptive header:

\begin{lstlisting}[language=Python]
header = ["OpenTime", "Open", "High", "Low", "Close",
          "Volume", "CloseTime", "QuoteVolume", "Trades",
          "TakerBuyBase", "TakerBuyQuote", "Ignore"]

with open("NEW_GOLD.csv", "w", newline="") as f:
    writer = csv.writer(f)
    writer.writerow(header)
    for k in klines:
        writer.writerow(k)
\end{lstlisting}

\subsubsection{4.1.3. Dataset Characteristics}
The resulting \texttt{NEW\_GOLD.csv} contains 1,000 rows plus a header, totalling approximately 95 KB. Key statistical properties of the Close price column:

\begin{center}
\begin{tabular}{|l|l|}
\hline
\textbf{Statistic} & \textbf{Value (approximate)}\\
\hline
Count & 1,000 daily records\\
Date range & May 2023 -- May 2026\\
Min Close & $\approx$ \$1,800 / oz\\
Max Close & $\approx$ \$3,400 / oz\\
Mean Close & $\approx$ \$2,300 / oz\\
Std. Dev. & $\approx$ \$380 / oz\\
\hline
\end{tabular}
\end{center}

The dataset captures the significant gold bull run of 2024--2026, providing a training distribution that includes both low-volatility consolidation periods and high-volatility rally periods.

\subsection{4.2. ML Training Module (\texttt{GoldPrice\_CLI})}

\subsubsection{4.2.1. Data Model: \texttt{ModelInput} and \texttt{ModelOutput}}
The ML.NET data model is defined using two simple C\# classes. The \texttt{[LoadColumn(index)]} attribute maps CSV column indices (0-based, accounting for the header row) to class properties:

\begin{lstlisting}
public class ModelInput
{
    [LoadColumn(1)] public float Open   { get; set; }
    [LoadColumn(2)] public float High   { get; set; }
    [LoadColumn(3)] public float Low    { get; set; }
    [LoadColumn(6)] public float Volume { get; set; }
    
    [LoadColumn(4), ColumnName("Label")]
    public float Close { get; set; }
}

public class ModelOutput
{
    [ColumnName("Score")]
    public float PredictedClose { get; set; }
}
\end{lstlisting}

Note that \texttt{Close} is decorated with both \texttt{[LoadColumn(4)]} (column index in the CSV) and \texttt{[ColumnName("Label")]} (the ML.NET convention for the regression target column). The \texttt{Volume} field is at index 6 (not index 5) because index 5 is the Close column in Binance's candlestick format, and index 6 corresponds to the actual traded volume.

\subsubsection{4.2.2. Data Loading and Splitting}
The program locates the CSV file using a two-step path resolution strategy --- parent directory first, then hardcoded fallback --- to handle both development (running from the project root) and containerised (running from a \texttt{/app/publish} directory) contexts:

\begin{lstlisting}
string projectRoot = Directory.GetParent(
    Directory.GetCurrentDirectory()).FullName;
DataPath = Path.Combine(projectRoot, "NEW_GOLD.csv");

if (!File.Exists(DataPath))
    DataPath = @"D:\Coding Space\Project\Group3\NEW_GOLD.csv";

var mlContext = new MLContext(seed: 42);
IDataView dataView = mlContext.Data.LoadFromTextFile<ModelInput>(
    path: DataPath, hasHeader: true, separatorChar: ',');

var split = mlContext.Data.TrainTestSplit(dataView, testFraction: 0.2);
\end{lstlisting}

\subsubsection{4.2.3. Training Pipeline Construction and Fitting}
The pipeline is constructed using ML.NET's fluent builder pattern. The \texttt{Concatenate} transform merges the four feature columns (which ML.NET requires to be explicitly listed) into a single \texttt{"Features"} vector column that the \texttt{LightGbm} trainer expects:

\begin{lstlisting}
var pipeline = mlContext.Transforms
    .Concatenate("Features",
        nameof(ModelInput.Open),
        nameof(ModelInput.High),
        nameof(ModelInput.Low),
        nameof(ModelInput.Volume))
    .Append(mlContext.Regression.Trainers.LightGbm(
        labelColumnName: "Label",
        featureColumnName: "Features",
        learningRate: 0.005,
        numberOfIterations: 1000));

Console.WriteLine("Training LightGBM model...");
ITransformer model = pipeline.Fit(split.TrainSet);
\end{lstlisting}

The call to \texttt{pipeline.Fit(split.TrainSet)} is synchronous and takes approximately 3--8 seconds on a modern CPU, depending on CPU speed and available cores. LightGBM internally parallelises histogram construction across CPU cores using OpenMP.

\subsubsection{4.2.4. Model Serialisation}
After training, the complete fitted pipeline (including the \texttt{Concatenate} transformer and the LightGBM model weights) is serialised to a portable ZIP archive:

\begin{lstlisting}
mlContext.Model.Save(model, split.TrainSet.Schema, "GoldModel.zip");
Console.WriteLine($"[OK] Model saved to {ModelPath}");
\end{lstlisting}

The \texttt{GoldModel.zip} file encapsulates:
\begin{itemize}
  \item The LightGBM model binary (native format with all tree structures and leaf values).
  \item The ML.NET schema definition mapping column names to data types.
  \item The \texttt{Concatenate} transformer configuration.
\end{itemize}

This self-contained artifact can be loaded on any .NET platform without retraining.

\subsubsection{4.2.5. Model Evaluation}
The trained model is applied to the 200-sample test set and evaluated using three regression metrics:

\begin{lstlisting}
var predictions = model.Transform(split.TestSet);
var metrics = mlContext.Regression.Evaluate(
    predictions,
    labelColumnName: "Label",
    scoreColumnName: "Score");

Console.WriteLine($"R2   : {metrics.RSquared:0.####}");
Console.WriteLine($"MAE  : {metrics.MeanAbsoluteError:0.####}");
Console.WriteLine($"RMSE : {metrics.RootMeanSquaredError:0.####}");
\end{lstlisting}

\subsubsection{4.2.6. Interactive Prediction REPL}
After evaluation, the CLI enters an interactive loop that prompts the user to enter OHLCV values manually:

\begin{lstlisting}
var engine = mlContext.Model.CreatePredictionEngine
             <ModelInput, ModelOutput>(model);
while (true) {
    Console.Write("> Enter Open price (or 'exit'): ");
    string input = Console.ReadLine();
    if (input?.ToLower() == "exit") break;
    float openVal = float.Parse(input);
    // ... read High, Low, Volume ...
    
    var sample = new ModelInput {
        Open = openVal, High = highVal,
        Low = lowVal, Volume = volVal
    };
    var result = engine.Predict(sample);
    Console.ForegroundColor = ConsoleColor.Green;
    Console.WriteLine(
        $"[RESULT] Predicted Close: {result.PredictedClose:0.00} USD");
    Console.ResetColor();
}
\end{lstlisting}

The REPL uses a regular (non-pooled) \texttt{PredictionEngine<T>} instance, which is acceptable in this single-threaded CLI context. The engine is created once before the loop and reused across all iterations, avoiding repeated model deserialisation overhead.

\subsection{4.3. API Server Module (\texttt{GoldPrice\_API})}

\subsubsection{4.3.1. Application Bootstrap (\texttt{Program.cs})}
The application entry point follows the clean \texttt{WebApplication} builder pattern introduced in .NET 6. All service registrations are delegated to extension methods to keep \texttt{Program.cs} concise and readable:

\begin{lstlisting}
var builder = WebApplication.CreateBuilder(args);

builder.Services.AddEndpointsApiExplorer();
builder.Services.AddSwaggerGen();
builder.Services.AddCors(options => options.AddPolicy("AllowAll",
    b => b.AllowAnyOrigin().AllowAnyMethod().AllowAnyHeader()));

builder.Services.AddCustomRateLimiting();  // Registers FixedWindow + HttpClient
builder.Services.AddGoldPriceModel();      // Registers PredictionEnginePool

var app = builder.Build();
app.UseSwagger();
app.UseSwaggerUI();
app.UseCors("AllowAll");
app.UseRateLimiter();
app.UseDefaultFiles();
app.UseStaticFiles();
app.MapPredictionEndpoints();
app.Run();
\end{lstlisting}

The middleware ordering is deliberate: Swagger before CORS (to allow Swagger UI to make cross-origin API calls), CORS before Rate Limiter (to attach CORS headers before rate limit rejection), static files after rate limiter (static assets are not rate-limited).

\subsubsection{4.3.2. Service Extensions (\texttt{ServiceExtensions.cs})}

\textbf{Rate Limiting Registration}:
\begin{lstlisting}
public static IServiceCollection AddCustomRateLimiting(
    this IServiceCollection services)
{
    services.AddRateLimiter(options => {
        options.AddFixedWindowLimiter("fixed", policy => {
            policy.PermitLimit = 5;
            policy.Window = TimeSpan.FromSeconds(10);
            policy.QueueProcessingOrder =
                QueueProcessingOrder.OldestFirst;
            policy.QueueLimit = 2;
        });
        options.RejectionStatusCode = 429;
    });
    services.AddHttpClient();
    return services;
}
\end{lstlisting}

\textbf{Model Pool Registration}:
\begin{lstlisting}
public static IServiceCollection AddGoldPriceModel(
    this IServiceCollection services)
{
    string modelPath = Path.Combine(
        Directory.GetCurrentDirectory(), "..",
        "GoldPrice_CLI", "GoldModel.zip");
    
    if (!File.Exists(modelPath))
        modelPath = @"D:\Coding Space\Project\Group3\
                      GoldPrice_CLI\GoldModel.zip";

    services.AddPredictionEnginePool<ModelInput, ModelOutput>()
        .FromFile(modelName: "GoldModel",
                  filePath: modelPath,
                  watchForChanges: true);
    return services;
}
\end{lstlisting}

\subsubsection{4.3.3. Realtime Prediction Endpoint}
The fully automated realtime endpoint demonstrates the highest technical complexity in the project, orchestrating three concurrent concerns (Binance API, Exchange Rate API, and ML inference) within a single asynchronous handler:

\begin{lstlisting}
apiGroup.MapGet("/realtime", async (
    PredictionEnginePool<ModelInput, ModelOutput> pool,
    IHttpClientFactory factory) =>
{
    var client = factory.CreateClient();
    float open=0, high=0, low=0, volume=0, vndRate=25000f;

    try {
        // 1. Fetch live OHLCV (today's candlestick)
        var res = await client.GetAsync(
            "https://api.binance.com/api/v3/klines"
            + "?symbol=PAXGUSDT&interval=1d&limit=1");
        res.EnsureSuccessStatusCode();
        var doc = JsonDocument.Parse(
            await res.Content.ReadAsStringAsync());
        var k = doc.RootElement[0];
        open   = float.Parse(k[1].GetString(),
                     CultureInfo.InvariantCulture);
        high   = float.Parse(k[2].GetString(),
                     CultureInfo.InvariantCulture);
        low    = float.Parse(k[3].GetString(),
                     CultureInfo.InvariantCulture);
        volume = float.Parse(k[5].GetString(),
                     CultureInfo.InvariantCulture);

        // 2. Fetch USD/VND exchange rate
        var er = await client.GetAsync(
            "https://open.er-api.com/v6/latest/USD");
        if (er.IsSuccessStatusCode) {
            var erDoc = JsonDocument.Parse(
                await er.Content.ReadAsStringAsync());
            vndRate = erDoc.RootElement
                .GetProperty("rates")
                .GetProperty("VND").GetSingle();
        }
    }
    catch { return Results.StatusCode(500); }

    // 3. Run ML inference
    var input = new ModelInput {
        Open=open, High=high, Low=low, Volume=volume };
    var pred = pool.Predict("GoldModel", input);

    return Results.Ok(new {
        success = true,
        data = new {
            inputsUsed = new { open, high, low, volume },
            predictedClose = pred.PredictedClose,
            predictedCloseVnd = pred.PredictedClose * vndRate,
            liveWorldPrice = open,
            exchangeRateVnd = vndRate
        }
    });
})
.WithName("GetRealtimeGoldPricePrediction")
.RequireRateLimiting("fixed");
\end{lstlisting}

\subsection{4.4. Web Interface (\texttt{wwwroot/index.html})}

\subsubsection{4.4.1. Design Language: Glassmorphism}
The X-AURUM web interface adopts the \textbf{Glassmorphism} design trend, characterised by:
\begin{itemize}
  \item \textbf{Frosted glass cards}: Achieved via \texttt{backdrop-filter: blur(12px)} and semi-transparent\\\texttt{background: rgba(255,255,255,0.08)} backgrounds.
  \item \textbf{Dark mode palette}: A deep navy/charcoal background gradient using\\\texttt{linear-gradient(135deg, \#0f0c29, \#302b63, \#24243e)} that conveys premium financial software aesthetics.
  \item \textbf{Accent colours}: Gold-toned highlights (\texttt{\#FFD700}) for price values and primary action elements, reinforcing the gold price theme.
  \item \textbf{Smooth transitions}: CSS \texttt{transition: all 0.3s ease} on interactive elements.
\end{itemize}

\subsubsection{4.4.2. Automatic Price Loading}
On page load, the JavaScript immediately calls the realtime endpoint without any user interaction, displaying the live predicted gold price within approximately 500 ms:

\begin{lstlisting}[language=Python]
// In index.html <script>
async function loadRealtimePrice() {
    const response = await fetch('/api/v1/predictions/realtime');
    const result = await response.json();
    if (result.success) {
        document.getElementById('predicted-usd').textContent
            = result.data.predictedClose.toFixed(2) + ' USD';
        document.getElementById('predicted-vnd').textContent
            = (result.data.predictedCloseVnd / 1e6).toFixed(2)
              + ' M VND';
        document.getElementById('live-price').textContent
            = result.data.liveWorldPrice.toFixed(2) + ' USD';
        document.getElementById('exchange-rate').textContent
            = result.data.exchangeRateVnd.toLocaleString() + ' VND/USD';
    }
}
window.addEventListener('DOMContentLoaded', loadRealtimePrice);
\end{lstlisting}

\subsubsection{4.4.3. Manual Prediction Form}
A secondary card on the page exposes a manual prediction form with four input fields (Open, High, Low, Volume). On form submission, the JavaScript POSTs the entered values to \texttt{/api/v1/predictions} and displays the result:

\begin{lstlisting}[language=Python]
document.getElementById('predict-form').addEventListener(
    'submit', async (e) => {
    e.preventDefault();
    const body = {
        open:   parseFloat(document.getElementById('inp-open').value),
        high:   parseFloat(document.getElementById('inp-high').value),
        low:    parseFloat(document.getElementById('inp-low').value),
        volume: parseFloat(document.getElementById('inp-volume').value)
    };
    const response = await fetch('/api/v1/predictions', {
        method: 'POST',
        headers: {'Content-Type': 'application/json'},
        body: JSON.stringify(body)
    });
    const result = await response.json();
    document.getElementById('manual-result').textContent
        = result.data.predictedClose.toFixed(2) + ' USD';
});
\end{lstlisting}

\subsubsection{4.4.4. Static File Serving}
The \texttt{index.html} file is placed in \texttt{wwwroot/} and served by ASP.NET Core's built-in static file middleware (\texttt{app.UseDefaultFiles(); app.UseStaticFiles()}). This eliminates the need for a separate web server (Nginx, Apache) to serve the frontend, keeping the Docker image architecture simple --- the single \texttt{ghcr.io/thanhtrunggdev/group3:latest} container serves both the API and the web UI.

\newpage
